\section{Differentiating the nonlinear insertion expansion}
\label{sec:leaf-expansion}
The elementary insertion is obtained by differentiating a real tilt of the
same nonlinear block integral used for the background. This matters for two
reasons. The background normalization must cancel exactly, and the derivative
of the gas must retain the original activity addresses. We first prove the
summation argument which makes this differentiation legitimate, then account
for the nonprincipal terms.

\subsection{A joint bound for values and derivatives}
Fix a number $D\ge0$ and an open interval $I$ containing zero. For a
differentiable real function $f$ and a nonnegative number $p$, say that
$f$ has joint bound $p$ at $\theta$ if there are $a,b\ge0$ with
\begin{equation}\label{eq:leaf-joint-bound}
 a+b\le p,\qquad |f(\theta)|\le a,\qquad
                         |f'(\theta)|\le Db.
\end{equation}
The witnesses $a,b$ may depend on $\theta$. This is stronger than
bounding the value by $p$ and the derivative by $Dp$ separately. Its
advantage is that products have no extra factor equal to their degree.
Indeed, if $f_1,f_2$ have witnesses $(a_1,b_1),(a_2,b_2)$, take
\[
       a=a_1a_2,\qquad b=a_1b_2+b_1a_2.
\]
The product rule gives the required value and derivative bounds, and
$a+b\le(a_1+b_1)(a_2+b_2)\le p_1p_2$. Induction gives joint
bound $\prod_{j\in q}p_j$ for any finite product. A constant multiplier
$c$ multiplies the bound by $|c|$. These statements include $D=0$.
For the affine tilt $f+\theta g$, if $|f|\le a$, $|g|\le Db$ and
$|\theta|D\le1$, the witnesses $(a+b,b)$ give joint bound $a+2b$.

Here is the precise infinite-sum consequence. Suppose $f_i$ is
differentiable on $I$ and has joint bound $p_i$ at every point of $I$,
where $\sum_i p_i<\infty$. Both $\sum_i f_i(\theta)$ and
$\sum_i f_i'(\theta)$ converge absolutely. For a closed segment
joining $\theta$ and $\theta+h$ inside $I$, the mean value theorem gives
\[
 \left|\frac{f_i(\theta+h)-f_i(\theta)}h\right|\le Dp_i.
\]
Dominated convergence on the counting measure therefore yields
\begin{equation}\label{eq:leaf-differentiate-sum}
             \frac{d}{d\theta}\sum_i f_i(\theta)
                         =\sum_i f_i'(\theta).
\end{equation}
This proof needs no continuity of the individual derivatives. It also
shows that the choice of an enumeration of the countable index set has
no effect on the result.

\subsection{The gas derivative and its support}
Work on a fixed finite coarse lattice $\Lambda$. Let $A_j\ne\varnothing$
be the support of forced activity $j\in\mathbb N$. The ordinary
activities, the support and inventory of each forced activity, and the
admissibility test are held fixed. Only the values $K_j$ change, along
a real family $F_j(\theta,U)$, with $F_j(0,U)=K_j(U)$.
Assume that, for $|\theta|D\le1$, these are measurable,
gauge invariant forced activities with the same locality and flag
conditions as $K_j$. Also assume
\begin{equation}\label{eq:leaf-season-bounds}
 \begin{split}
 &F_j(\theta,\cdot)=F_j(0,\cdot)\quad\hbox{if }o\notin A_j,\\
 &F_j(\cdot,U)\text{ has joint bound }p_j\ge0,\\
 &\sup_{x\in\Lambda}\sum_{j:x\in A_j}p_j e^{u C_B(A_j)}\le s,
                         \qquad u\ge0.
 \end{split}
\end{equation}
The middle line is required on an open neighborhood of zero contained
in $\{|\theta|D\le1\}$, for every $U$. Such a neighborhood exists
also when $D=0$. Since each support is nonempty,
\[
                  \sum_jp_j\le |\Lambda|s<\infty.
\]
For a finite family $q$ put $p_q=\prod_{j\in q}p_j$.
Expanding finite products and taking increasing finite subsets of
$\mathbb N$ proves
\[
       \sum_{\substack{q\subset\mathbb N\;\mathrm{finite}}}p_q
              =\prod_j(1+p_j)\le\exp\Bigl(\sum_jp_j\Bigr)<\infty.
\]
For fixed $U$, admissibility does not depend on $\theta$. Thus each
admissible product $\prod_{j\in q}F_j(\theta,U)$ has joint bound
$p_q$, and an inadmissible term is zero. Equation~\eqref{eq:leaf-differentiate-sum}
now proves that the gas derivative at zero is
\begin{equation}\label{eq:leaf-gas-derivative}
 \mathcal G^{[1]}(k)(U)
 =\sum_{q\ {\rm admissible}}\sum_{j\in q}
             k_j(U)\prod_{\ell\in q\setminus\{j\}}K_\ell(U),
 \qquad k_j(U)=\left.\partial_\theta F_j(\theta,U)\right|_0.
\end{equation}
There is also absolute convergence before the marked index is summed.
For joint-bound witnesses $a_j,b_j$ at zero,
\[
 \sum_{j\in q}|F_j'(0,U)|\prod_{\ell\ne j}|K_\ell(U)|
 \le D\sum_{j\in q}b_j\prod_{\ell\ne j}a_\ell
 \le D\prod_{j\in q}(a_j+b_j)\le Dp_q.
\]
The middle inequality keeps just the terms with exactly one $b$ in
the product expansion. Summing over $q$ proves absolute convergence
of the full pointed series. This is the marked gas on the original addresses.

The marked family retains all required structural properties.
If $o\notin A_j$, its derivative is zero by the first line of
\eqref{eq:leaf-season-bounds}. Wherever the flag condition requires
$F_j(\theta,U)=0$, it also requires $k_j(U)=0$. Gauge invariance and
locality follow by differentiating the corresponding equalities in
$\theta$. Measurability follows by choosing any nonzero sequence
$\theta_n\to0$ in the permitted interval and writing $k_j$ as the
pointwise limit of measurable difference quotients. Finally,
\[
 |k_j(U)|\le Dp_j\,\mathbf1_{\{o\in A_j\}},\qquad
 \sum_j Dp_j\mathbf1_{\{o\in A_j\}}e^{uC_B(A_j)}\le Ds.
\]
This proves the flagged cost \eqref{eq:leaf-flag-input}, with exactly
the insertion debt $D$ and the site bound $s$ in
\eqref{eq:leaf-season-bounds}. Complexification preserves the same
majorant.

\subsection{The logarithmic derivative and the exact raw identity}
The localized regular terms arise from the logarithm of the positive
normalization. Write this normalization as
\[
 a(\theta,U)=\exp\left(\sum_i\ell_i(\theta,U)\right),
                 \qquad a_0(U)=a(0,U)>0.
\]
Suppose that each logarithmic term has a summable joint majorant,
uniformly for $\theta$ near zero, and that its derivative is the
localized insertion
\[
               \ell_i'(0,U)=\chi_{X_i}(U)h_i(U).
\]
The product bound above applies, in particular, to each connected
cluster term $c\prod_j f_j(\theta)/r!$: its majorant is
$|c|\prod_jp_j/r!$. Thus an absolutely summable connected-cluster
majorant justifies differentiating the logarithmic expansion.
Equation~\eqref{eq:leaf-differentiate-sum} and the chain rule give
\begin{equation}\label{eq:leaf-log-derivative}
      a'(0,U)=a_0(U)z(U),\qquad
                 z(U)=\sum_i\chi_{X_i}(U)h_i(U).
\end{equation}

The exact tilted block identity, after the two activity regroupings,
has the form
\begin{equation}\label{eq:leaf-tilted-identity}
 \mathcal T\rho_{\rm fine}(U)+\theta\mathcal T f(U)
       =e^{E(U)}\frac{a(\theta,U)}{a_0(U)}\,b(\theta,U),
 \qquad b(0,U)=\mathcal G(K)(U),
\end{equation}
where $b$ is the gas of $F(\theta)$. All quantities use the same
background and history. This identity is required for the actual block
integral, for $|\theta|D\le1$. Differentiating it, using
\eqref{eq:leaf-gas-derivative} and \eqref{eq:leaf-log-derivative}, yields
\[
 \mathcal T f(U)
 =e^{E(U)}\mathcal G(K)(U)z(U)
                      +e^{E(U)}\mathcal G^{[1]}(k)(U).
\]
As $\rho_{\rm coarse}=e^E\mathcal G(K)$, this is exactly
\eqref{eq:leaf-raw-identity}. In particular, the derivative of the
normalization is multiplied by the full coarse density. Omitting
$a_0^{-1}$ before differentiation would give an incorrect factor.

The normalization cancellation in \eqref{eq:leaf-tilted-identity}
is algebraic, but uses exact identities. To spell it out, let $G_\theta$
be the gas before the two regroupings, whose logarithmic shifts are
$R$ and $S$. Their identities read
\[
 b(\theta)=G_\theta e^{R+S},\qquad
 e^{E+R+S}=h a_0,\qquad
 h a(\theta)G_\theta=\mathcal T\rho_{\rm fine}+\theta\mathcal T f.
\]
Substitution proves \eqref{eq:leaf-tilted-identity}. The two exact gas regroupings and their numerical majorants are
inputs here, as is the exact block expansion in the last identity.
Their application must retain the same support, inventory and flag
data. The scalar differentiation argument does not
by itself establish the additional transferred singleton-profile
identity required in Section~\ref{sec:leaf-transfer}.

\subsection{Separating the principal term from the tail}
Here is the quantitative bookkeeping for the nonprincipal regular
terms. It identifies precisely which connected estimates are needed.
Let $j$ be the derivative order, $p$ the spatial weight exponent,
$B$ the cover radius, and $a\ge0$, $e>0$ the support exponent and its
reserve. Put
\[
 r=pj,\qquad \Phi(C)=[\sigma(1+C)]^r,\qquad
 P_{e,r}=e^{-r}r!\exp(e),\qquad
                          A_{e,r}=\sigma^r P_{e,r},
\]
where $\sigma=4(2B+1)>0$ and $e^{-r}$ denotes the inverse power of
the reserve $e$. The exponential series gives
$[e(1+C)]^r/r!\le\exp(e(1+C))$ for $C\ge0$. Consequently
\begin{equation}\label{eq:leaf-polynomial-reserve}
                 \Phi(C)e^{aC}\le A_{e,r}e^{(a+e)C}.
\end{equation}

The localized logarithmic terms have a nonnegative integer order.
An order-zero term corresponds to a single thin molecule; the higher
orders are connected clusters. Terms which do not contain the insertion
vanish. The remaining terms have supports containing the chosen
anchor $o$. Write $C_i$ for their cover numbers and $\omega_i\ge0$
for their cluster majorants. Suppose the connected estimate gives,
for every $z\ge0$ and every site $x$,
\begin{equation}\label{eq:leaf-cluster-tail}
 \sum_{i:\,x\in X_i,\,\operatorname{ord}(i)\ge z}
              \omega_i e^{(a+e)C_i}
                      \le t\gamma^z(1-\gamma)^{-1},
                        \qquad 0\le\gamma<1.
\end{equation}
For positive-order terms the derivative estimate is
$\|h_i\|\le D\omega_i\Phi(C_i)$. Keeping only terms containing the
insertion, using $o\in X_i$, and applying
\eqref{eq:leaf-polynomial-reserve} and
\eqref{eq:leaf-cluster-tail} at $z=1$ gives the total bound
\begin{equation}\label{eq:leaf-positive-tail}
 \sum_{\operatorname{ord}(i)\ge1}e^{aC_i}\|h_i\|
                   \le D A_{e,r}\,t\gamma(1-\gamma)^{-1}.
\end{equation}

Order zero must be treated separately, because its single-insertion
molecule is the principal term. If an order-zero molecule containing
the insertion is not principal, it has at least two elements. Indeed,
a singleton containing the marked element consists exactly of that
element. Thus its excess cardinality is at least one. Let $v_i$ be
its nonnegative marked-molecule majorant, with the restriction that
excess cardinality is at least one. The mixed marked estimate needed
here is
\[
       \sup_x\sum_{i:\,x\in X_i}v_i e^{(a+e)C_i}\le C_1.
\]
Assume the order-zero derivative bound
$\|h_i\|\le D B_*v_i\Phi(C_i)$ for a nonnegative derivative price
$B_*$. The molecule-to-order-zero-cluster map is injective, so no
multiplicity is lost in the sum. Applying the reserve inequality and
the site estimate at $o$ proves
\[
 \sum_{\substack{\operatorname{ord}(i)=0\\i\ {\rm nonprincipal}}}
                  e^{aC_i}\|h_i\|\le D B_* A_{e,r} C_1.
\]
Together with \eqref{eq:leaf-positive-tail}, this gives
\begin{equation}\label{eq:leaf-tail-produced}
 M\le D A_{e,r}\left(B_*C_1+\frac{t\gamma}{1-\gamma}\right).
\end{equation}
To obtain the reserved mass of Section~\ref{sec:leaf-transfer}, take
$a$ at least its reserved support exponent. The same calculation with
a larger $a$ continues to apply if the displayed estimates have that
reserve.

In the actual insertion normalization,
$D=m\exp(-\varepsilon C_{\rm fine})/g$, where $g>0$ is the gain and
$\varepsilon>0$ is the fine support reserve. Therefore, defining
\[
 T_{\rm tail}=\frac{A_{e,r}}g
                   \left(B_*C_1+\frac{t\gamma}{1-\gamma}\right),
\]
we have $M\le m e^{-\varepsilon C_{\rm fine}}T_{\rm tail}$.
This accounts for both the support decay and the normalization gain;
they are not independent assumptions on $M$.

\subsection{The strict power reserve}
For completeness, the tail comparison used in choosing the coupling
can be verified directly from the price exponents. Write
$H(x)=\beta_{\rm fine}(x)$, with $H(x)\to\infty$. A function has
upper exponent $q$ if for every $v>q$, eventually
$|f(x)|\le H(x)^v$. Products add upper exponents and sums with the
same upper exponent retain it: use half the available exponent margin
in each bound and absorb the fixed factor two into the remaining
power of $H$. Fixed powers of a function with upper exponent zero
still have upper exponent zero.

Let $q<0$ be the connected-expansion exponent. The prices used above
have the following exponents, with $L$, derivative orders and $e$
fixed before $x$ is taken large:
\[
\begin{array}{c|cccccc}
\text{price}&A_{e,r}&B_*&C_1&t&\gamma&g^{-1}\\ \hline
\text{upper exponent}&0&0&3q/2&q&q/2&-q.
\end{array}
\]
Here $\gamma=H^{q/2}$ and $g=H^q$; hence eventually
$0\le\gamma\le1/2$ and $(1-\gamma)^{-1}\le2$. Both terms defining
$T_{\rm tail}$ consequently have upper exponent $q/2$.
Since $q/2<q/4$, the definition gives, eventually,
\begin{equation}\label{eq:leaf-tail-power}
                 T_{\rm tail}\le H^{q/4},\qquad
          M\le m e^{-\varepsilon C_{\rm fine}}H^{q/4}.
\end{equation}
This is a strict exponent comparison. An upper exponent zero by itself
would not justify paying a price from a fixed power of $\log H$.

Sections~\ref{sec:cluster-expansion} and~\ref{sec:mixed-activity-prices}
prove the connected and mixed site estimates from their vertex prices;
Section~\ref{sec:soft-log-fusion} proves the two gas regrouping identities
under their exact flag and reserve hypotheses. The remaining analytic
work is to construct the actual tilted block expansion, establish its
vertex prices and localized derivative estimates, and purchase all
these bounds and their stated upper exponents on the selected history.
With those inputs, the differentiation, flagged cost, nonprincipal
tail accounting and strict-power comparison above supply the
corresponding inputs to the elementary projection estimate.
